\documentclass[exercice,a4]{thedocument}%
\usepackage[french]{babel}%
\usepackage{textcomp}%
\usepackage{pstricks}%
\usepackage{multido}%
\usepackage{scratch3}%
\usepackage{tabularx}%
\startdatakeys
    \source{21GENMATPO1}
    \BO{2026}
    \cycle{4}
    \competencesDuSocle{CA,RA}
    \objectifsApprentissage{C4D2ba,C4A1af}
    \motsCles{Scratch,programme,variable}
\enddatakeys
\begin{document}%
\def\degres{\textsuperscript{\textdegree}}%

\medskip

On donne le programme suivant:

\medskip

\begin{tabularx}{\linewidth}{m{9cm} m{0.25cm} X}
\begin{scratch}
\setscratch{num blocks=true}
\blockinit{quand \greenflag est cliqué}
\blockmove{aller à  x : \ovalnum{190} y : \ovalnum{0}}
\blockmove{s'orienter à \ovalnum{90} \degres}
\blockvariable{mettre{\selectmenu{Longueur} à \ovalnum{30}}}
\blockrepeat{répéter \ovalnum{4} fois}
{\blocklook{Motif}
\blockpen{relever le stylo}
\blockmove{avancer de \ovalmoreblocks{Longueur} * \ovalnum{2} + \ovalnum{10}}
\blockmove{ajouter à \selectmenu{Longueur} \ovalnum{10}}}
\end{scratch}&
\renewcommand{\arraystretch}{2}
\begin{tabular}{l}
~\\
A\\
B\\
C\\
D\\
E\\
\end{tabular}&
\renewcommand{\arraystretch}{1}
\begin{scratch}
%\newcommand*\numblock[1]{\color{black}\footnotesize\bfseries#1}
\initmoreblocks{définir \namemoreblocks{Motif}}
\blockpen{stylo en position d'écriture}
\blockrepeat{répéter \ovalnum{6} fois}
{\blockmove{avancer de \ovalmoreblocks{Longueur}}
\blockmove{tourner \turnright{} de \ovalnum{60} degres}
}
\end{scratch}
\end{tabularx}

\medskip

On rappelle que \og s'orienter à 90 \fg{} signifie que l'on est orienté vers la droite.

\medskip

\begin{enumerate}
\item On prendra dans cette question $1$~mm pour un pixel.

Représenter en vraie grandeur sur votre copie la figure que trace le bloc Motif lorsque Longueur
vaut $30$ pixels.
\item Ce programme utilise une variable, quel est son nom ? À quoi correspond-elle sur la figure
réalisée par le bloc Motif ?
\item Laquelle de ces trois figures obtient-on lorsqu'on exécute ce programme ? Indiquer sur la copie le numéro de la bonne proposition parmi les trois suivantes. On expliquera son choix.

\begin{center}
\begin{tabularx}{\linewidth}{|p{2.25cm}|X|X|}\hline
\psset{unit=1cm}
\begin{pspicture}(2,-2)
\rput(1,-0.2){1}
\def\motif{\pspolygon(0.4;0)(0.4;60)(0.4;120)(0.4;180)(0.4;240)(0.4;300)}
\rput(0.4,-0.8){\motif}
\rput(0.46,-0.86){\psscalebox{1.2 1.2}{\motif}}
\rput(0.53,-0.93){\psscalebox{1.44 1.44}{\motif}}
\rput(0.6,-1){\psscalebox{1.728 1.728}{\motif}}
\end{pspicture}&\begin{pspicture}(6.5,-2)
\rput(2.5,-0.2){2}
\def\motif{\pspolygon(0.4;0)(0.4;60)(0.4;120)(0.4;180)(0.4;240)(0.4;300)}
\rput(0.4,-0.9){\motif}\rput(1.5,-0.98){\pspolygon(0.48;0)(0.48;60)(0.48;120)(0.48;180)(0.48;240)(0.48;300)}\rput(2.75,-1.076){\pspolygon(0.576;0)(0.576;60)(0.576;120)(0.576;180)(0.576;240)(0.576;300)}\rput(4.2,-1.1912){\pspolygon(0.6912;0)(0.6912;60)(0.6912;120)(0.6912;180)(0.6912;240)(0.6912;300)}
\end{pspicture}&\begin{pspicture}(6.5,-2)
\rput(2.5,-0.2){3}
\rput(0.5,-1.1912){\pspolygon(0.6912;0)(0.6912;60)(0.6912;120)(0.6912;180)(0.6912;240)(0.6912;300)}\rput(2,-1.1){\pspolygon(0.576;0)(0.576;60)(0.576;120)(0.576;180)(0.576;240)(0.576;300)}
\rput(3.3,-1){\pspolygon(0.48;0)(0.48;60)(0.48;120)(0.48;180)(0.48;240)(0.48;300)}
\rput(4.5,-0.95){\pspolygon(0.4;0)(0.4;60)(0.4;120)(0.4;180)(0.4;240)(0.4;300)}
\end{pspicture}\\ \hline
\end{tabularx}
\end{center}
\item Modifier le programme précédent pour obtenir la figure ci-dessous. Pour cela, indiquer les numéros des instructions à supprimer ou à modifier, et préciser les modifications à apporter:

\begin{center}
\psset{unit=1cm}
\begin{pspicture}(0,-0.5)(9.5,0.5)
%\psgrid
\def\motif{\pspolygon(0.4;0)(0.4;60)(0.4;120)(0.4;180)(0.4;240)(0.4;300)}
\multido{\n=2.00+1.25}{6}{\rput(\n,0){\motif}}
\end{pspicture}
\end{center}

\item On souhaite modifier le bloc Motif afin qu'il permette de tracer un carré. Pour cela, indiquer les lettres des instructions à supprimer ou à modifier, et préciser les modifications à apporter.
\end{enumerate}

\bigskip

\showthesource
\end{document}
